\documentclass[11pt]{article}

% ---------- Packages ----------
\usepackage[margin=1in]{geometry}
\usepackage{amsmath,amssymb}
\usepackage{booktabs}
\usepackage{array}
\usepackage{enumitem}
\usepackage{xcolor}
\usepackage[most]{tcolorbox}
\usepackage{titlesec}
\usepackage{parskip}
\usepackage{hyperref}
\usepackage{fancyhdr}

% ---------- Hyperref ----------
\hypersetup{
  colorlinks=true,
  linkcolor=blue!50!black,
  urlcolor=blue!50!black,
  pdftitle={AP Statistics: AI Tutoring, Compliance, and Evidence (Student Edition)}
}

% ---------- Header / footer ----------
\pagestyle{fancy}
\fancyhf{}
\fancyhead[L]{\small AP Statistics --- Student Edition}
\fancyhead[R]{\small \thepage}
\renewcommand{\headrulewidth}{0.4pt}

% ---------- Section formatting ----------
\titleformat{\part}[display]
  {\normalfont\Huge\bfseries\centering}
  {\partname\ \thepart}{0.5em}{}
\titleformat{\section}
  {\normalfont\Large\bfseries\color{blue!50!black}}{\thesection}{1em}{}
\titleformat{\subsection}
  {\normalfont\large\bfseries\color{blue!40!black}}{\thesubsection}{1em}{}

% ---------- Boxes ----------
\newtcolorbox{scenario}{
  colback=gray!5, colframe=gray!50,
  title=Scenario, fonttitle=\bfseries,
  boxrule=0.5pt, arc=2pt
}
\newtcolorbox{notebox}{
  colback=blue!3, colframe=blue!40!black,
  title=Note, fonttitle=\bfseries,
  boxrule=0.5pt, arc=2pt
}

% ---------- Custom commands ----------
\newcommand{\code}[1]{\texttt{#1}}
\newcommand{\phat}{\hat{p}}
\newcommand{\xbar}{\bar{x}}

\title{\bfseries AP Statistics:\\ AI Tutoring, Compliance, and Evidence \\[0.4em]
\large A Combined Capstone and Investigative Task \\[0.2em]
\normalsize Student Edition}
\author{}
\date{}

\begin{document}
\maketitle
\thispagestyle{empty}

\begin{center}
\begin{minipage}{0.85\textwidth}
\itshape
This document combines two complementary AP Statistics investigations.
\textbf{Part I} is a 60-question synthetic capstone covering Units 1--9.
\textbf{Part II} is a focused investigative task that applies those same
tools to a real-world question about AI-tutor pedagogy. Show all work.
Cite the question number for each answer. Use $\alpha = 0.05$ unless
otherwise stated.
\end{minipage}
\end{center}

\vspace{1em}
\tableofcontents
\newpage

% ============================================================
\part{Capstone: AI Tutoring, Compliance, and Evidence}
% ============================================================

\begin{notebox}
This is a synthetic AP Statistics investigation. The scenario is
realistic, but the data are invented so the statistical reasoning can
be evaluated cleanly.
\end{notebox}

\section{Scenario}

A district is debating whether AI tutoring should be allowed during the
learning phase of AP Statistics. Some teachers argue that AI mainly
enables cheating. Others argue that AI works like a private tutor:
useful during practice, but not a substitute for independent reasoning
on assessments.

The district runs a six-week study with 240 AP Statistics students
from 8 schools. Within each school, students are randomly assigned to
one of two instructional conditions:

\begin{itemize}[leftmargin=*]
  \item \textbf{AI-access condition:} Students may use an AI tutor
  during practice, homework, and review. They must still show
  handwritten work and complete in-class checks without AI.
  \item \textbf{Standard-resource condition:} Students use AP
  Classroom videos, teacher notes, textbook examples, and ordinary
  peer discussion. They may not use AI during the study period.
\end{itemize}

Researchers measure each student's pretest score, posttest score,
daily AI-tutor minutes, phone-policy environment, engagement rating,
and whether the student passes a later AP-style practice exam. Use a
significance level of $\alpha = 0.05$ unless otherwise stated.

\section{Data}

\subsection*{Table A: One-Variable Data}
For 20 randomly selected students in the AI-access condition, the
recorded average daily AI-tutor minutes were:
\[
0,\ 0,\ 5,\ 8,\ 10,\ 12,\ 15,\ 18,\ 22,\ 25,\ 30,\ 32,\ 35,\ 40,\ 45,\ 50,\ 55,\ 60,\ 75,\ 95
\]

\subsection*{Table B: Randomized Study Summary}
\begin{center}
\small
\begin{tabular}{lrrrrrrrr}
\toprule
Condition & $n$ & Mean pre & SD pre & Mean post & SD post & Mean gain & SD gain & Passed \\
\midrule
AI-access          & 120 & 48.2 & 13.8 & 66.8 & 15.6 & 18.6 & 12.4 & 69 \\
Standard-resource  & 120 & 48.9 & 14.1 & 58.7 & 16.2 &  9.8 & 11.7 & 47 \\
\bottomrule
\end{tabular}
\end{center}

\subsection*{Table C: Actual AI Use Within the AI-Access Condition}
\begin{center}
\small
\begin{tabular}{lrrrr}
\toprule
Actual behavior & $n$ & Mean gain & SD gain & Passed \\
\midrule
Regular use ($\geq 20$ min/day)        & 76 & 23.4 & 11.6 & 55 \\
Little or no use ($< 20$ min/day)      & 44 & 10.3 & 10.9 & 14 \\
\bottomrule
\end{tabular}
\end{center}

\subsection*{Table D: Separate Observational Survey of School Culture and AI Use}
The following table comes from a separate districtwide observational
survey. It is not part of the randomized experiment in Table B.
\begin{center}
\small
\begin{tabular}{lrrrr}
\toprule
School culture & Regular & Occasional & None & Total \\
\midrule
Restrictive phone/AI culture & 21 & 39 & 60 & 120 \\
Open AI-tutor culture        & 66 & 34 & 20 & 120 \\
\midrule
\textbf{Total}               & \textbf{87} & \textbf{73} & \textbf{80} & \textbf{240} \\
\bottomrule
\end{tabular}
\end{center}

\subsection*{Table E: Regression Summary}
For 80 students with complete usage logs, researchers fit a
least-squares regression model predicting gain score from average
daily AI-tutor minutes:
\[
\widehat{\text{gain}} = 7.0 + 0.40(\text{minutes})
\]
Additional information:
\begin{itemize}[leftmargin=*,nosep]
  \item $n = 80$
  \item $r = 0.62$
  \item $s = 9.8$
  \item $SE_b = 0.061$ (standard error of the slope)
  \item Residual plot shows random scatter with roughly constant vertical spread.
  \item Histogram of residuals is roughly unimodal and approximately symmetric.
\end{itemize}

\subsection*{Table F: Probability Model}
Based on a larger district survey, suppose an AP Statistics student
independently has probability $p = 0.32$ of becoming a regular
AI-tutor user when AI access is allowed but not required.

\section{Questions}

\subsection*{Unit 1: Exploring One-Variable Data}
\begin{enumerate}[leftmargin=*]
\item Classify each variable as categorical or quantitative:
instructional condition, school culture, average daily AI-tutor
minutes, gain score, passed AP-style practice exam.
\item Use Table A to make an appropriate graph of AI-tutor minutes.
Describe shape, center, spread, and unusual features in context.
\item Find the median, $Q_1$, $Q_3$, IQR, and range of Table A.
\item Use the $1.5 \times \text{IQR}$ rule to determine whether 95
minutes is an outlier.
\item Explain why the mean AI-tutor minutes would be larger than the
median for Table A.
\end{enumerate}

\subsection*{Unit 2: Exploring Two-Variable Data}
\begin{enumerate}[resume,leftmargin=*]
\item Use Table D to compute the conditional relative frequency of
regular AI use within each school culture.
\item Does Table D suggest an association between school culture and
actual AI use? Explain using conditional relative frequencies.
\item Interpret the slope $0.40$ in Table E in context.
\item Interpret $r = 0.62$ in context.
\item Find and interpret $r^2$.
\item Predict the gain score for a student who uses the AI tutor for
40 minutes per day.
\item Explain why it would be risky to use the regression line to
predict the gain score for a student who uses AI for 180 minutes
per day.
\end{enumerate}

\subsection*{Unit 3: Collecting Data}
\begin{enumerate}[resume,leftmargin=*]
\item Identify the experimental units, treatments, and response
variables in the six-week study.
\item Explain why random assignment within each school improves the
study.
\item Explain why blocking by school would be sensible in this study.
\item Can the randomized comparison in Table B support a
cause-and-effect conclusion about AI access? Explain.
\item Can the comparison in Table C support a cause-and-effect
conclusion about actual AI use? Explain carefully.
\item Identify one possible source of bias if students self-report
their AI minutes.
\item Explain how noncompliance can contaminate the interpretation of
a treatment effect.
\end{enumerate}

\subsection*{Unit 4: Probability, Random Variables, and Probability Distributions}
\begin{enumerate}[resume,leftmargin=*]
\item Using Table D, find $P(\text{regular AI use})$.
\item Using Table D, find $P(\text{open AI-tutor culture})$.
\item Using Table D, find $P(\text{regular AI use} \mid \text{open AI-tutor culture})$.
\item Are regular AI use and open AI-tutor culture independent in
Table D? Justify numerically.
\item Suppose 5 students are randomly selected from a population
where $p = 0.32$ for regular AI-tutor use. Find the probability
exactly 2 are regular AI-tutor users.
\item Let $X$ be the number of regular AI-tutor users among 5
randomly selected students with $p = 0.32$. Find $E(X)$ and
$\text{SD}(X)$.
\item Find the probability that the first regular AI-tutor user
appears on the fourth student selected.
\end{enumerate}

\subsection*{Unit 5: Sampling Distributions}
\begin{enumerate}[resume,leftmargin=*]
\item For samples of size 120 from a population with true pass
proportion $p = 0.50$, describe the approximate sampling
distribution of $\hat{p}$. Check conditions.
\item If the true mean gain under the standard-resource condition is
10 points with population SD 12, describe the approximate sampling
distribution of the sample mean gain for $n = 120$.
\item Under the null hypothesis that the two pass proportions are
equal, explain why the observed difference in pass rates from
Table B might or might not be surprising.
\item Explain why larger sample sizes make it easier to detect
smaller real differences.
\end{enumerate}

\subsection*{Unit 6: Inference for Categorical Data: Proportions}
\begin{enumerate}[resume,leftmargin=*]
\item Construct a 95\% confidence interval for the proportion of
AI-access students who passed the AP-style practice exam.
\item Construct a 95\% confidence interval for the difference in pass
proportions, $p_{AI} - p_{\text{standard}}$, using Table B.
\item Interpret the interval from Question 32 in context.
\item Test whether the AI-access condition produced a higher pass
proportion than the standard-resource condition.
\item Explain the meaning of the p-value from Question 34 in context.
\item Describe a Type I error and a Type II error for the test in
Question 34.
\end{enumerate}

\subsection*{Unit 7: Inference for Quantitative Data: Means}
\begin{enumerate}[resume,leftmargin=*]
\item Construct a 95\% confidence interval for the mean gain in the
AI-access condition.
\item Construct a 95\% confidence interval for the difference in mean
gains, $\mu_{AI} - \mu_{\text{standard}}$, using Table B.
\item Interpret the interval from Question 38 in context.
\item Test whether the AI-access condition produced a greater mean
gain than the standard-resource condition.
\item Explain why the randomized design matters when interpreting the
difference in mean gains.
\item Suppose a critic says, ``Students in the AI-access group were
already stronger.'' Use Table B to respond statistically.
\end{enumerate}

\subsection*{Unit 8: Inference for Categorical Data: Chi-Square}
\begin{enumerate}[resume,leftmargin=*]
\item Use Table D to state appropriate hypotheses for a chi-square test.
\item Calculate the expected count for restrictive culture and
regular AI use.
\item Verify that the conditions for a chi-square test are met.
\item The chi-square test statistic for Table D is approximately
$43.62$. Find the degrees of freedom and approximate p-value
category.
\item State a conclusion in context.
\item Identify which cells contribute most to the chi-square
statistic and explain what that means in context.
\end{enumerate}

\subsection*{Unit 9: Inference for Quantitative Data: Slopes}
\begin{enumerate}[resume,leftmargin=*]
\item State hypotheses for testing whether average daily AI-tutor
minutes are linearly associated with gain score.
\item Use Table E to calculate the test statistic for the slope.
\item Find the degrees of freedom.
\item State a conclusion in context.
\item Construct a 95\% confidence interval for the slope.
\item Interpret the interval in context.
\item Explain why regression evidence alone is weaker causal evidence
than the randomized comparison in Table B.
\end{enumerate}

\subsection*{Final Synthesis}
\begin{enumerate}[resume,leftmargin=*]
\item A school administrator says, ``The AP-style scores prove
AI-supported instruction works.'' Explain why that statement is
too strong.
\item A teacher says, ``The AP-style scores prove AI-supported
instruction fails whenever students do not use it correctly.''
Explain why that statement is also too strong.
\item Distinguish between the effect of being assigned to the
AI-access condition and the association observed for students who
actually used AI regularly.
\item Write a statistically careful conclusion about what the full
dataset does and does not show.
\item Propose a better follow-up study that could distinguish
AI-tutor effects from student motivation, compliance, teacher
culture, and phone-policy effects.
\end{enumerate}

\newpage
% ============================================================
\part{Investigative Task: Evaluating an AI-Tutor Pedagogy}
% ============================================================

\begin{scenario}
A high school AP Statistics teacher in Massachusetts implemented an
AI-tutor-supplemented curriculum during the 2024--25 school year. The
class had $n = 22$ students. At year's end, \textbf{all 22 scored a 1}
on the AP exam. The national 1-rate that year was $\mathbf{23.7\%}$
(the 17.0\% figure was the national 5-rate). The
teacher reports that 18 of the 22 students ``boycotted'' the exam
(left it largely blank as protest); the remaining 4 attempted it.

Separately, a published randomized study (Kestin et al., 2025) compared
a scaffolded AI tutor to research-based active-learning instruction on
identical content. The study reported substantially higher post-test
performance in the AI-tutor condition ($p < 10^{-8}$).

The teacher's father, citing the all-1s outcome, argues
\emph{``the data shows your method failed.''}
\end{scenario}

\section{(a) Sampling and external validity \textnormal{\small(Unit 3)}}

The teacher's class is \textbf{not} a random sample from any larger
population.
\begin{enumerate}[leftmargin=*,label=(\roman*)]
  \item Identify the population the local result can validly describe.
  \item Identify the population the published study can validly describe.
  \item Explain why averaging the two as if equally weighted commits a
  \textbf{selection bias} error.
\end{enumerate}

\section{(b) Inference on the local outcome \textnormal{\small(Units 4 \& 6)}}

State appropriate hypotheses for testing whether the local 1-rate
differs from the national 1-rate ($p = 0.237$). Check the conditions
for a one-proportion $z$-test. If the conditions fail, propose and
carry out an appropriate alternative test, compute the p-value, and
state a conclusion in context.

\section{(c) Noncompliance and the meaning of the outcome \textnormal{\small(Units 3 \& 6)}}

Of the 22 students, 18 deliberately did not attempt the exam.
Explain in context what the all-1s outcome actually measures, and
identify what data would be more informative about the causal effect
of the AI-tutor curriculum on statistical reasoning skill.

\section{(d) Interpreting the published study \textnormal{\small(Units 3 \& 7)}}

The published randomized study reported substantially higher post-test
performance in the AI-tutor condition than in the active-learning
condition ($p < 10^{-8}$). State what causal claim is justified by
this evidence and what the scope of inference is. Be precise about the
comparator (research-based active learning, not traditional lecture).

\section{(e) ``We don't have the data yet'' claim \textnormal{\small(Units 3 \& 6)}}

A critic argues, ``We don't have the data on AI tutoring.'' Use the
published study to evaluate this claim. Be careful about what is and
is not justified by a single study.

\section{(f) Confounding in the local sample \textnormal{\small(Unit 3)}}

List at least three plausible confounding variables in the local
outcome that are not present in the published randomized study.
Explain why randomization addresses these confounders in the
published study but not in the local classroom outcome.

\section{(f2) Design matters: tutoring versus crutching \textnormal{\small(Unit 3)}}

Bastani et al.\ (2025, \emph{PNAS}) studied nearly 1,000 high-school
math students. During practice, unrestricted GPT-4 access improved
assisted performance. But when AI was removed on later exams, those
students performed \emph{worse} than the control group. A scaffolded
tutor condition (hints rather than full answers) largely eliminated
the harm.

Explain what this means for the question ``Should AI be allowed in
classrooms?'' and why ``AI access'' cannot be treated as a single
treatment. Identify the two distinct treatments hidden inside ``AI
access.''

\section{(g) Chi-square test for differential response \textnormal{\small(Unit 8)}}

The teacher records the following $2 \times 2$ table from the next
cohort ($n = 30$), classifying students by prior engagement and
post-tutor performance:

\begin{center}
\small
\begin{tabular}{lrrr}
\toprule
 & Improved $\geq$ 1 letter grade & Did not improve & Total \\
\midrule
Previously disengaged & 9 & 3  & 12 \\
Previously engaged    & 4 & 14 & 18 \\
\midrule
\textbf{Total}        & \textbf{13} & \textbf{17} & \textbf{30} \\
\bottomrule
\end{tabular}
\end{center}

State appropriate hypotheses, check conditions, compute the
chi-square statistic, find the degrees of freedom and an approximate
p-value, and state a conclusion in context.

\section{(h) Linear regression with interaction \textnormal{\small(Units 2 \& 9)}}

Researchers fit:
\[
\widehat{\text{Score}} = \beta_0 + \beta_1(\text{AI hours})
  + \beta_2(\text{prior engagement})
  + \beta_3(\text{AI hours} \times \text{prior engagement})
\]
The fitted model gives $\hat\beta_1 > 0$ and $\hat\beta_3 < 0$.
Interpret what these signs together imply about how the marginal
benefit of AI hours depends on a student's prior engagement.

\section{(i) Synthesis: stating the defensible conclusion \textnormal{\small(Skills 1.E, 4.B)}}

Write a paragraph, in AP-rubric style, that:
\begin{enumerate}[leftmargin=*]
  \item Identifies the \textbf{causal question}.
  \item Identifies \textbf{which evidence is admissible} for that
  question and which is not.
  \item States the \textbf{direction} of the admissible evidence.
  \item Identifies the \textbf{failure mode} of the local outcome.
  \item States a defensible empirical claim, with appropriate scope
  of inference, and explains why arguments based only on the local
  all-1s outcome rely on inferential errors.
\end{enumerate}

\end{document}
